% !TeX root = ../main.tex
\section{Roofline Interpretation}\label{sec:roofline}
\rubric{Roofline and quantitative insight}{20}{Correct arithmetic intensity, use of measured memory
bandwidth, and a justified compute- vs.\ memory-bound conclusion.}

\begin{hypothesis}
  All four kernels have very low arithmetic intensity (well under \qty{1}{\flop\per\byte}), so on a
  machine whose single-core bandwidth is tens of \unit{\giga\byte\per\second} and whose peak vector
  rate is over \qty{100}{\giga\flop\per\second} they should all sit on the \emph{bandwidth} slope of
  the roofline when DRAM-resident, i.e.\ be memory-bound. The stencil, with the highest intensity,
  should come closest to the compute roof.
\end{hypothesis}
\begin{hypothesis}
  (For the measured ceilings.) Even the L1 bandwidth roof should put the
  ridge point above \qty{0.5}{\flop\per\byte}, so the cache-resident native kernels should also be
  bounded by bandwidth (or by the load/store ports that define it), not by the FMA units.
\end{hypothesis}

Attainable performance at arithmetic intensity $I$ (FLOPs per byte of memory traffic) is bounded by
\begin{equation}
  P(I) = \min\bigl(P_{\text{peak}},\; B_{\ell} \cdot I\bigr),
\end{equation}
where $P_{\text{peak}}$ is the core's peak FLOP rate and $B_{\ell}$ the bandwidth of the memory level
$\ell$ the data lives in \cite{williams2009roofline}. We draw one bandwidth roof per level (a
cache-aware roofline). Both ceilings are \emph{measured by separate microbenchmarks}
(\cref{sec:ceilings-method}), never derived from the kernels plotted against them.

\begin{table}[tbp]
  \centering\small
  \caption{Measured roofline ceilings, single core (\code{run.sh ceilings}). Top: register-only
  accumulator loops (\num{12} independent chains); FLOP/cycle and clock from the user-space cycle
  counter. Bottom: STREAM-style bandwidth at one footprint per level (GiB/s). \enquote{incl.\ WA}
  adds the write-allocate read; the roofs in \cref{fig:roofline} use the explicit triad in L1
  (stores hit L1, no write-allocate) and triad incl.\ WA below it. The read stream ($^{*}$) is
  add-latency-bound (\cref{sec:ceilings-method}) and is not used as a roof.}\label{tab:ceilings}
  \IfDataExists{ceilings_peak}{%
    \pgfplotstabletypeset[
      col sep=comma,
      columns={isa,gflops_median,flops_per_cycle,ghz},
      columns/isa/.style={column name=Peak loop, string type, string replace*={_}{\_}},
      columns/gflops_median/.style={column name=GFLOP/s, fixed, zerofill, precision=1},
      columns/flops_per_cycle/.style={column name=FLOP/cycle, fixed, zerofill, precision=2},
      columns/ghz/.style={column name=GHz, fixed, zerofill, precision=2},
      every head row/.style={before row=\toprule, after row=\midrule},
      every last row/.style={after row=\bottomrule},
    ]{data/ceilings_peak.csv}%
  }\par\medskip
  \IfDataExists{ceilings_bw}{%
    \pgfplotstabletypeset[
      col sep=comma,
      columns={level,bytes,read,copy,triad,triad_total},
      columns/level/.style={column name=Level, string type},
      columns/bytes/.style={column name=Footprint (B), int detect},
      columns/read/.style={column name=read$^{*}$, fixed, zerofill, precision=1},
      columns/copy/.style={column name=copy, fixed, zerofill, precision=1},
      columns/triad/.style={column name=triad, fixed, zerofill, precision=1},
      columns/triad_total/.style={column name={triad incl.\ WA}, fixed, zerofill, precision=1},
      every head row/.style={before row={\toprule & & \multicolumn{4}{c}{GiB/s}\\}, after row=\midrule},
      every last row/.style={after row=\bottomrule},
    ]{data/ceilings_bw.csv}%
  }
\end{table}

\begin{table}[tbp]
  \centering\small
  \caption{Arithmetic intensity (float32) from the actual traffic per element, the same convention
  as the triad roof: L2, L3 and DRAM include the write-allocate read of a store to an unread line;
  in L1 every store hits, so there is none.}\label{tab:roofline-inputs}
  \begin{tabularx}{\linewidth}{@{}lrrX@{}}
    \toprule
    Kernel & $I$, L2--DRAM & $I$, L1 & Traffic per element (L2--DRAM) \\
    \midrule
    SAXPY        & \num{0.167} & \num{0.167} & $2\,\text{FLOP} / \qty{12}{\byte}$: read $x$, read $y$, write $y$
                                  ($y$ is read first, so no write-allocate) \\
    multiply     & \num{0.063} & \num{0.083} & $1\,\text{FLOP} / \qty{16}{\byte}$: read $x,y$, write $z$ + write-allocate of $z$ \\
    stencil      & \num{0.417} & \num{0.625} & $5\,\text{FLOP} / \qty{12}{\byte}$: read $x$ once (reuse), write $y$ + write-allocate \\
    dot product  & \num{0.25}  & \num{0.25}  & $2\,\text{FLOP} / \qty{8}{\byte}$: read $x,y$, no store \\
    triad (roof) & \num{0.125} & \num{0.167} & $2\,\text{FLOP} / \qty{16}{\byte}$: read $b,c$, write $a$ + write-allocate \\
    \bottomrule
  \end{tabularx}
\end{table}

\begin{figure}[tbp]
  \centering
  \IfDataExists{roofline_points}{%
    \begin{tikzpicture}
      \begin{axis}[acs, xmode=log, ymode=log, xlabel={Arithmetic intensity (FLOP/byte)},
                   ylabel={Performance (GFLOP/s)}, width=0.7\linewidth, height=0.6\linewidth,
                   xmin=0.02, xmax=20, ymin=0.5, ymax=400,
                   legend pos=outer north east, legend style={font=\scriptsize}]
        \IfFileExists{data/roofline_roofs.csv}{%
          \addplot[thick, black] table[x=ai, y=gflops, only if={roof}{L1}] {data/roofline_roofs.csv};
          \addlegendentry{L1 roof}
          \addplot[thick, black!70, densely dashed] table[x=ai, y=gflops, only if={roof}{L2}] {data/roofline_roofs.csv};
          \addlegendentry{L2 roof}
          \addplot[thick, black!55, dashdotted] table[x=ai, y=gflops, only if={roof}{L3}] {data/roofline_roofs.csv};
          \addlegendentry{L3 roof}
          \addplot[thick, black!40, dotted] table[x=ai, y=gflops, only if={roof}{DRAM}] {data/roofline_roofs.csv};
          \addlegendentry{DRAM roof}
          \addplot[thin, red!60] table[x=ai, y=gflops, only if={roof}{peak_sse_f32}] {data/roofline_roofs.csv};
          \addlegendentry{SSE peak (autovec)}
          \addplot[thin, red!60, densely dashed] table[x=ai, y=gflops, only if={roof}{peak_scalar_f32}] {data/roofline_roofs.csv};
          \addlegendentry{scalar peak}
        }{}
        \addplot+[only marks, mark=*, nodes near coords, point meta=explicit symbolic,
                  every node near coord/.style={font=\tiny, anchor=south west}, iqr bars]
          table[x=ai, y=gflops_median, iqr of=gflops, meta=kernel, only if={variant}{native}, only if={level}{L1}] {data/roofline_points.csv};
        \addlegendentry{native, L1-resident}
        \addplot+[only marks, mark=square*, iqr bars] table[x=ai, y=gflops_median, iqr of=gflops, only if={variant}{native}, only if={level}{L2}] {data/roofline_points.csv};
        \addlegendentry{native, L2-resident}
        \addplot+[only marks, mark=triangle*, iqr bars] table[x=ai, y=gflops_median, iqr of=gflops, only if={variant}{native}, only if={level}{L3}] {data/roofline_points.csv};
        \addlegendentry{native, L3-resident}
        \addplot+[only marks, mark=diamond*, iqr bars] table[x=ai, y=gflops_median, iqr of=gflops, only if={variant}{native}, only if={level}{DRAM}] {data/roofline_points.csv};
        \addlegendentry{native, DRAM}
        \addplot+[only marks, mark=o, black, iqr bars] table[x=ai, y=gflops_median, iqr of=gflops, only if={variant}{scalar}, only if={level}{DRAM}] {data/roofline_points.csv};
        \addlegendentry{scalar, DRAM}
      \end{axis}
    \end{tikzpicture}%
  }
  \caption{Cache-aware roofline for one core (L1 points at their L1 intensity,
  \cref{tab:roofline-inputs}). Roofs: measured triad bandwidth per level, capped by the
  measured AVX2-FMA float32 peak; thin red lines: the measured SSE and scalar peaks that bound the
  \code{autovec} and \code{scalar} builds. Points: kernel GFLOP/s from \cref{sec:baseline} at
  $N=\num{1024}$, \num{16384}, \num{524288} and \num{16777216} (L1, L2, L3, DRAM), median with IQR error bars (mostly smaller than the markers).}\label{fig:roofline}
\end{figure}

\subsection{Analysis}
\emph{Ceilings} (\cref{tab:ceilings}). The AVX2 FMA loop reaches \qty{89.2}{\giga\flop\per\second}
float32 (\num{44.6} float64), which is \num{31.99}~FLOP/cycle at the measured
\qty{2.79}{\giga\hertz}. That is exactly the $2~\text{FMA ports} \times 8~\text{lanes} \times 2 = 32$
FLOP/cycle the core can do, so the ports, not the loop, set this ceiling. The SSE mul+add loop
(\code{autovec}'s ISA, no FMA) reaches \qty{22.3}{\giga\flop\per\second} (\num{8.00}~FLOP/cycle,
expected $2\times4=8$) and the scalar loop \qty{5.55}{\giga\flop\per\second}
(\num{1.99}~FLOP/cycle, expected $2$). Triad bandwidth (actual traffic) falls from
\qty{155.6}{\gibi\byte\per\second} in L1 to \num{90.5}, \num{48.6} and
\qty{24.3}{\gibi\byte\per\second} in L2, L3 and DRAM. In L1 this is the explicit load and store
traffic: the triad's stores hit lines already in L1, so the write-allocate read counted below L1
(\qty{207}{\gibi\byte\per\second} incl.\ WA) does not happen there. The ridge points are
$I = P_{\text{peak}}/B_\ell \approx \num{0.53}$, \num{0.92}, \num{1.71} and
\qty{3.42}{\flop\per\byte}. The L1 ridge is just above the \qty{0.5}{\flop\per\byte} the second
hypothesis assumed. Only the stencil, whose L1 traffic is \qty{8}{\byte} per element
($I=\num{0.625}$), lies right of it, i.e.\ nominally under the compute roof in L1.

\emph{DRAM-resident kernels.} Every kernel's intensity (\numrange{0.06}{0.42}~FLOP/byte) is far left
of the DRAM ridge, so all four are bandwidth-bound in DRAM. This is consistent with the hypothesis
and with the locality sweep (\cref{sec:locality}). Against the DRAM roof $B_{\text{DRAM}}\cdot I$,
native multiply reaches \SI{100}{\percent} (\num{1.63} measured vs \qty{1.63}{\giga\flop\per\second}),
the stencil \SI{101}{\percent} and SAXPY \SI{107}{\percent} (\num{4.65} vs \num{4.35}). Multiply has
exactly the triad's access pattern (two reads, one write to an unread line) and lands on the roof.
SAXPY, whose pattern differs (two reads, one write to a line it just read), exceeds it by
\SI{7}{\percent}, and by up to \SI{20}{\percent} at the intermediate $N=\num{8388608}$
(\cref{sec:negative}). The sustainable bandwidth depends a little on the access pattern, so the
triad roof is a reference, not an upper bound to the byte. Because the roof comes from a different
loop, a kernel landing on it is evidence, not a tautology. The scalar SAXPY, multiply and stencil
DRAM points sit at \SIrange{37}{66}{\percent} of the same roof (scalar dot: \SI{20}{\percent},
latency-bound like native dot). They are the largely instruction-bound loops of
\cref{sec:locality}, which is why SIMD still pays $1.51\times$
(multiply), $1.94\times$ (SAXPY) and $2.72\times$ (stencil) at these points.

\emph{Cache-resident kernels.} The L3-resident native points sit at \SIrange{95}{100}{\percent} of the
L3 roof (SAXPY \SI{100}{\percent}, multiply \SI{100}{\percent}, stencil \SI{95}{\percent}), so L3
bandwidth bounds them exactly. Closer to the core they fall below the roof. In L2, SAXPY reaches
\SI{86}{\percent} and multiply \SI{88}{\percent} of the roof (stencil \SI{42}{\percent}). In L1,
native SAXPY delivers \qty{20.1}{\giga\flop\per\second} at $I=\num{0.167}$, \SI{72}{\percent} of the
L1 roof and \SI{23}{\percent} of the FMA peak; multiply reaches \SI{66}{\percent} of the L1 roof and
the stencil only \SI{19}{\percent} of the peak that bounds it there. So cache-resident kernels are
bounded neither by the FMA units (confirming the second hypothesis) nor quite by the measured L1
bandwidth. The limit is the loop's own load/store and instruction issue:
\numrange{0.22}{0.28}~cycles per element for SAXPY versus the \num{0.125} that two loads and a store
per cycle would allow (\cref{sec:baseline}). The roofline thus explains the step pattern of
\cref{fig:locality} from L2 outward, since each step is a move down to a lower roof. Inside L1/L2 it
only gives the upper bound.

\emph{The dot product} is the one kernel far \emph{below} every roof:
\qtyrange{1.33}{1.49}{\giga\flop\per\second} at all sizes (\cref{sec:baseline}). It is not
bandwidth-bound but latency-bound by the non-reassociated reduction (\cref{sec:vectorization}). One
scalar \code{vaddss} per element with a \num{4}-cycle latency gives an in-core ceiling of
$2~\text{FLOP}/4~\text{cycles} = 0.5$~FLOP/cycle, i.e.\ \qty{1.39}{\giga\flop\per\second} at the
measured \qty{2.79}{\giga\hertz}. The measured \qtyrange{1.39}{1.42}{\giga\flop\per\second} for every
$N$ from \num{1024} to \num{1048576} matches it to within \SI{2}{\percent}; only $N\le512$ runs slightly
above it (\qtyrange{1.45}{1.49}{\giga\flop\per\second}). Each call starts a new, independent chain,
and with short chains the out-of-order core overlaps the end of one call with the next. DRAM runs
slightly below it (\qty{1.33}{\giga\flop\per\second}). The roofline, a throughput model, does not apply to the dot
product until the dependency chain is broken. To move a streaming kernel toward the compute roof one
would have to raise its arithmetic intensity (fuse passes, or block for cache so data is re-read from
L1 rather than DRAM); widening the vectors alone cannot help a kernel that is already waiting on
memory.
